
\lesson{23/11}
So the quantities to be computed are probability amplitudes of the form
\[
    A=\phantom{x}_\text{out}\bra{q_1...q_m}\ket{p_1...p_n}_\text{in}
\]
where $p_i,q_j$ are momenta eigenstate.

For non-free theories it was found that
\[
    ia^\dagger(t,\mathbf{p})=\int\dd^3x g_p\bar{\pd_t} \phi(t,\mathbf{x})
\]
\[
    -ia(t,\mathbf{p})=\int\dd^3x g^*_p\bar{\pd_t} \phi(t,\mathbf{x})
\]
Their derivative was found to be
\[
    i\pd_ta^\dagger (t,\mathbf{p}) = \int_{}^{} \dd^3x g_p(\Box+m^2)\phi(t,
    \mathbf{x})
\]

A strategy to compute the amplitudes uses
\begin{gather*}
    \lim\limits_{T\to-\infty} a^{(\dagger)}(t,p) = \sqrt{Z}a_\text{in}^{(\dagger)}
    (p)\\
    \lim\limits_{T\to+\infty} a^{(\dagger)}(t,p) =\sqrt{Z}a_\text{out}^{(\dagger)}
    (p),
\end{gather*}
since
\[
    a_\text{in}^\dagger (\mathbf{p})\ket{0} = \sqrt{2E_p}\ket{p},
\]
the amplitude can be written as
\[
    \phantom{x}_\text{out}\bra{q_1...q_m}\frac{1}{\sqrt{2E_p}}a^\dagger_\text{in}
    (p_1)\ket{p_2...p_n}_\text{in}
\]
adding $0=a_\text{out}^\dagger(p_1)-a_\text{out}^\dagger(p_1)$ yields
\[
    \frac{1}{\sqrt{2E_p}}\phantom{x}_\text{out}\bra{q_1...q_m}a_\text{in}^\dagger
    -a_\text{out}^\dagger(p_1)\ket{p_2...p_n}_\text{in} +
    \frac{1}{\sqrt{2E_p}}\phantom{x}_\text{out}\bra{q_1...q_m}a^\dagger_\text{out}
    (p_1)\ket{p_2...p_n}_\text{in}.
\]
Assuming \textit{there is no forward scattering}, i.e. that no out particle has
the same momentum as any in particle (which means no particle remains
``untouched'' by the interaction), implies that the second addendum is zero (
since the $a^\dagger $ commutes with all bra creation operators that
create the out state and acts on the $\bra{0}$, yielding 0):
\[
    A = \frac{1}{\sqrt{2E_p}}\frac{-1}{\sqrt{Z}}\lim\limits_{T\to\infty}
    \phantom{x}_\text{out}
    \bra{q_1...q_m}
    \left(a^\dagger(T,p_1)-a^\dagger(-T,p_1)\right)\ket{p_2...p_n}_\text{in}
\]
The difference of the operators can be written as the integral over time of their
derivative
\begin{gather*}
    A = \frac{1}{\sqrt{2E_p}}\frac{-1}{\sqrt{Z}}\int_{-\infty}^{\infty} \dd t
    \phantom{x}_
    \text{out}\bra{q_1...q_m}\pd_ta^\dagger(t,p_1)\ket{p_2...p_n}_\text{in}=\\
    =\frac{1}{\sqrt{2E_p}}\frac{i}{\sqrt{Z}}\int_{}^{} \dd^4x_1g_{p_1}(x_1)
    (\Box_1+m^2)\phantom{x}_\text{out}\bra{q_1...q_m}\phi(x_1)\ket{p_2...p_n}_
    \text{in}
\end{gather*}
This process can be iterated to express $A$ in terms of the vacuum expectation
value of time ordered products of fields: one can 
express
\[
    \ket{p_2...p_n}_\text{in} = a^\dagger_\text{in}(p_2)\ket{p_3...p_n}
\]
and add
\[
    0=\left(a^\dagger_\text{out}(p_2)-a^\dagger_\text{out}(p_2)\right)\phi(x_1)
\]
and repeat the same process as before:
\[
    -\frac{1}{\sqrt{Z}}\lim\limits_{t_2\to\infty}\bra{q_1...q_m}a^\dagger(
    t_2,p_2)\phi(x_1)-a^\dagger(-t_2,p_2)\phi(x_1)\ket{p_3...p_n}\propto
    \int_{-\infty}^{\infty} \dd t_2 \pd_{t_2}\bra{q_1...q_m}T(a^\dagger(t_2,p_2)
    \phi(x_1)\ket{p_3...p_n}
\]  
and expressing $\pd_{t_2}a^\dagger(t_2,p_2)$ yields a term
\[
    \bra{q_1...q_m}T(\phi(x_2)\phi(x_1))\ket{p_3...p_n}.
\]
A similar strategy can be used to express $\bra{q_1...q_m}$ using $a$'s instead of
$a^\dagger$'s: after removing the initial states
\[
    \bra{q_1...q_m}T(\phi(x_1)...\phi(x_n))\ket{0} = 
    \bra{q_2...q_m}a_\text{out}(q_2)T(\phi_1...\phi_n) + T(...)(a_\text{in}(q_2)-
    a_\text{in}(q_2))\ket{0}
\]
from here, one can take the total time derivative, integrated over all time.

The final formula is
\begin{gather*}
A = \phantom{x}_\text{out}\bra{q_1..q_m}\ket{p_1..p_n}_\text{in}=\\=
\frac{i^n(-i)^m}{Z^{\frac{1}{2}(n+m)}}\int \left(\prod_{i=1}^{n} \dd^4x_i
\right)\left(\prod_{j=1}^{m}\dd^4y_j\right)\left(\prod_{i=1}^{n}\frac{g_{p_i}
(x_i)}{\sqrt{2E_{p_i}}}\right)\left(\prod_{j=1}^{m}\frac{g^*_{q_j}
(y_j)}{\sqrt{2E_{q_j}}}\right)(\Box_{x_1}+m^2)...(\Box_{x_n}+m^2)\\
(\Box_{y_1}+m^2)...(\Box_{y_m}+m^2)\bra{\Omega}T(\phi(y_1)...\phi(y_m)\phi(x_1)
...\phi(x_n))\ket{\Omega}.
\end{gather*}
The $g_p^{(*)}$ are plane waves and $x_i,y_j$ are integrated upon: these are
nothing but \textit{Fourier Transforms}.

The $\Box_{i}$'s become $-p_i^2$'s in momentum space:
\[
    A = \mathcal{N}\int\left(\prod_{i=1}^{n}\frac{\dd^4x_i}{(2\pi)^{\frac{3}{2}}}
    e^{ip_ix_i}\right)\left(\prod_{j=1}^{m}\frac{\dd^4y_j}{(2\pi)^{\frac{3}{2}}}
    e^{-iq_jy_j}\right)(-p_1^2+m^2)...(-q^2_m+m^2)\bra{\Omega}T(\phi_1...\phi_n...
    \phi_m)\ket{\Omega}
\]
\textbf{Note}: the particles are on shell, so the $(-p^2_k+m^2)$ should be all
zeroes, \textbf{but} the propagator has poles in the same points, so the
integrand is \textit{the amputated propagator}, since the poles 
associated to the external legs are simplified.

So at the end, the probability amplitude is
\textbf{the Fourier transform of the amputated propagator}:
\[
    A = \frac{i^n(-i)^m}{Z^{(n+m)/2}}\tilde{G}^\text{amp}(p_1...p_n;q_1...q_m)
    (2\pi)^4\delta^4\left(\sum p_i+\sum_{}^{}q_j\right)
\]

In principle, particles can be different and have different masses: they yield
different $Z$ factors, different mass terms  
\[
    A = \mathcal{N}\int\left(\prod_{i=1}^{n}\frac{\dd^4x_i}{(2\pi)^{\frac{3}{2}}}
    e^{ip_ix_i}\right)\left(\prod_{j=1}^{m}\frac{\dd^4y_j}{(2\pi)^{\frac{3}{2}}}
    e^{-iq_jy_j}\right)(-p_1^2+m_1^2)...(-q^2_m+m_m^2)\bra{\Omega}
    T(\phi_1(x_1)...\phi_n(x_n)...\phi_m(y_m))\ket{\Omega}
\]
where
\[
    \mathcal{N} = \frac{i^n(-i)^m}{\prod_{ij}Z_i^mZ_j^n}
\]
and $\phi_1...\phi_n...\phi_m$ are the $n+m$ different fields associated to the
$n+m$ different particles.
\subsection{Examples}
Let $\ket{in} = \ket{p_1,p_2}$ and $\ket{out} = \ket{q_1,q_2}$ in the
$\phi^4$ theory: there are two particles in and two particles out, so
the order is $\Lambda^2$ and there are 3 possible diagrams (since the
hypotesis of no forward scattering implies only connected diagrams):\\
\begin{tikzpicture}
\begin{feynman}
  \draw [scalar] (4, 0.5) arc [start angle=90, delta angle=180, radius=1];
  \draw [scalar] (4, 0.5) arc [start angle=90, delta angle=-180, radius=1];
  \vertex (v0) at (-2.5, 0.5) {};
  \vertex[dot] (v1) at (-1.5, -0.5) {};
  \vertex (v2) at (-2.5, -1.5) {};
  \vertex[dot] (v3) at (0.5, -0.5) {};
  \vertex (v4) at (1.5, 0.5) {};
  \vertex (v5) at (1.5, -1.5) {};
  \vertex (v6) at (2.5, 0.5) {};
  \vertex (v7) at (5.5, 0.5) {};
  \vertex (v8) at (2.5, -1.5) {};
  \vertex (v9) at (5.5, -1.5) {};
  \vertex[dot] (v10) at (4, 0.5) {};
  \vertex[dot] (v11) at (4, -1.5) {};
  \vertex (v12) at (6.5, 0.5) {};
  \vertex[empty dot, minimum size=0.025mm] (v13) at (8.5, 0.5) {};
  \vertex (v14) at (6.5, -1.5) {};
  \vertex[empty dot, minimum size=0.1mm] (v15) at (8.5, -1.5) {};
  \vertex[dot] (v16) at (7.5, 0.5) {};
  \vertex[dot] (v17) at (7.5, -1.5) {};
  \vertex (v18) at (10.5, -1.5) {};
  \vertex (v19) at (10.5, 0.5) {};
  \diagram* {
    (v0) -- [scalar] (v1),
    (v1) -- [scalar] (v2),
    (v1) -- [scalar, bend right=45] (v3),
    (v1) -- [scalar, bend left=45] (v3),
    (v3) -- [scalar] (v4),
    (v3) -- [scalar] (v5),
    (v6) -- [scalar] (v7),
    (v8) -- [scalar] (v9),
    (v12) -- [scalar] (v13),
    (v14) -- [scalar] (v15),
    (v16) -- [scalar, bend right=45] (v17),
    (v16) -- [scalar, bend left=45] (v17),
    (v15) -- [scalar] (v19)
  };
  \begin{scope}[even odd rule]
    \clip (-5, -4) rectangle (13, 3) (9.5, -0.5) circle (0.082cm);
    \path (v13) edge [scalar] (v18);
  \end{scope}
\end{feynman}
\end{tikzpicture}



Using Feynman rules the propagator can be computed as the sum of $3$ propagators


